\section{个人总结}
本人于2018年完成清华大学土木工程本科专业的学习，获得清华本科生特等奖学金（本科生最高荣誉），2020 年获得美国麻省理工学院(MIT)交通科学与计算机科学双硕士学位，2022 年获得麻省理工学院交通科学博士学位并获得 MIT UPS 博士奖学金。主要研究方向为城市公共交通系统，以及与该系统相关的需求与行为建模，政策分析，和多模式交通系统耦合。在研究中使用到多种数学模型，包括优化（鲁棒优化、整数规划、非线性优化），概率与统计、机器学习（深度学习、强化学习）、计量经济学与博弈论。在交通领域顶级期刊（TR-A, B, C, E, IEEE ITS, TS 等）共发表了\textbf{{\themyCounterNum}}篇学术论文（全为SCI，\textbf{\topnum}篇为中科院一区Top，总影响因子\totalif），其中\textbf{{\themyFirstAuthorNum}} 篇一作（一作总影响因子\firstauthorif），Google Scholar总引用量\textbf{\totalcitation}。于硕士在读期间开发了一套基于地铁刷卡数据的自适应网络服务状态监控系统，被应用于港铁的每日运营评估，服务于日均500万人次的日常出行。博士期间致力于城市公共交通系统的韧性研究，建立起了一套高效的公共交通事故影响评估系统，可以迅速分析和可视化事故影响，提出的乘客路径推荐模型可为芝加哥公共交通系统在事故期间降低20\% 的延误时间。研究曾两次被 MIT News（学校官方媒体）报道。本人也与工业界合作紧密，曾在 Amazon 世界范围内的路径规划挑战赛中获得全球第二名（唯一学生获奖队伍，第一名为三位滑铁卢大 学数学系终身教授），获得5万美元奖金。曾在美国第二大的共享出行公司Lyft担任高级研究科学家，负责公司最重要的实时司机激励产品“Bonus Map”的算法设计与迭代，提出的新模型为公司带来每年3800万美元收益（经过A/B Testing验证）。目前在TikTok（美国最大短视频平台）的电商团队担任高级机器学习科学家，负责平台美国地区物流补贴算法从0到1的开发, 将平台物流补贴回报率提升16.7\%, 销售额不变的情况下每天节约30万美元（经过A/B Testing验证）。2024年入选国家级青年人才项目（海外）。
% 